\chapter{Prakata}

Buku ini adalah buku referensi tentang arsitektur perangkat lunak. Isinya
membahas empat belas gaya arsitektur, dari client-server sampai
\textit{serverless}, dan setiap gaya dinilai dengan program yang benar-benar
dijalankan. Angka pada tabelnya adalah keluaran program tersebut. Angka itu
dikutip apa adanya, termasuk angka yang melawan harapan. Cara itulah yang
dimaksud kata terukur pada judulnya, dan seksi Arti Kata Terukur pada
Pendahuluan memberi contohnya.

Buku ini juga dapat digunakan sebagai buku ajar. Materinya berkembang dari modul
mata kuliah Arsitektur Perangkat Lunak (IF231303) di Universitas Pradita,
sehingga setiap bab membawa Tujuan Pembelajaran, \textit{Summary}, dan Latihan
beserta lembar isiannya. Bab~15 berisi proyek akhir untuk kelas. Bagian
Penggunaan sebagai Buku Ajar pada Pendahuluan merinci susunan tersebut.

Pembaca sasarannya ada dua kelompok. Kelompok pertama adalah pengembang dan
arsitek yang harus memilih gaya untuk sistem nyata, lalu ingin mengetahui harga
setiap pilihan sebelum keputusan diambil. Kelompok kedua adalah pengajar dan
mahasiswa yang membutuhkan bahan yang dapat dijalankan ulang, bukan sekadar
dibaca.

Seluruh program, skrip pengukuran, slide, dan sumber diagram tersedia di
\url{https://github.com/alfa-yohannis/software-architecture}. Kode tiap bab
berada di \berkas{sources/chapterNN/} dan slidenya di
\berkas{slides/sessionNN/}, dengan NN sama dengan nomor babnya. Setiap direktori
kode punya \berkas{README.md} yang mendaftar \textit{file}-nya beserta latihan
yang memakainya.

Terima kasih kepada mahasiswa mata kuliah Arsitektur Perangkat Lunak di
Universitas Pradita yang memakai bahan awal buku ini. Koreksi dari pembaca dapat
disampaikan lewat halaman \textit{issue} repositori tersebut atau lewat
penerbit.

\bigskip
\begin{flushright}
  Tangerang, \isian{bulan dan tahun terbit}\\[4pt]
  Penulis
\end{flushright}
